% ============================================================
%  FCFS Appointment Simulation Results
%  Supervisor brief
% ============================================================

\documentclass[11pt]{article}

\usepackage[margin=0.85in]{geometry}
\usepackage{amsmath}
\usepackage{booktabs}
\usepackage{tabularx}
\usepackage{array}
\usepackage{microtype}
\usepackage{parskip}
\usepackage{enumitem}
\usepackage[hidelinks]{hyperref}

\setlist[itemize]{leftmargin=1.25em,itemsep=2pt,topsep=3pt}
\setlist[enumerate]{leftmargin=1.5em,itemsep=3pt,topsep=3pt}
\newcolumntype{Y}{>{\raggedright\arraybackslash}X}
\newcolumntype{L}[1]{>{\raggedright\arraybackslash}p{#1}}

\title{FCFS Appointment Simulation Results\\[0.25em]
\large Supervisor Brief}
\author{Alexandre Sepulveda de Dietrich\\[0.2em]
\small Under the supervision of Prof.\ Carri Chan}
\date{May 2026}

\begin{document}
\maketitle

\section*{Executive Summary}

The main result is that capacity utilization and patient access separate sharply
in the current FCFS stress-test simulation. At baseline, average utilization is
\(0.839\), but only \(0.269\) of arrivals complete a visit. This is not a
contradiction. Utilization is normalized by the number of appointment slots,
while served rate is normalized by arrivals. With \(100\) arrivals per day and
\(32\) slots per day, the served rate is mechanically bounded near
\(32/100=0.32\), even before no-shows, cancellations, and balking.

The baseline loss pattern is dominated by balking and cancellation, not by a
lack of slot offers. Of all arrivals, \(26.9\%\) are served, \(35.6\%\) balk,
\(32.3\%\) cancel, \(5.2\%\) no-show, and \(0.0\%\) receive no offer. The zero
no-offer rate matters: at the baseline demand level, the horizon does not fully
saturate because daily cancellations keep freeing future capacity. Access
losses are therefore mainly behavioral/post-booking losses under long offered
delays, not immediate ``no slot available'' failures.

A controlled no-balking diagnostic shows that eliminating balking does not
increase completed utilization. Across 30 seeds, utilization changes from
\(0.8418\) to \(0.8404\), while booked rate rises by \(5.85\) percentage
points. The additional accepted bookings are offset by no-offer accumulation
\((0.0000 \rightarrow 0.2966)\) and higher cancellation exposure
\((0.3248 \rightarrow 0.3836)\). Under the current formulation, balking mainly
changes where unmet demand is recorded; it does not solve the completed-visit
bottleneck.

The sensitivity analysis gives a clear hierarchy of mechanisms. Total arrival
rate is the dominant driver of served rate and offered wait. No-show behavior
is the most direct driver of completed utilization because no-show slots are
not rebooked. Balking affects who receives access through the shared FCFS pool,
but has little direct effect on aggregate utilization near the baseline.
Accepted wait is a selective metric because high-delay patients are more likely
to balk; offered wait is a better congestion measure, but it must be reported
with no-offer share because patients with no offer are absent from all wait
metrics.

\textbf{Bottom line.} For this stress-test baseline, the most defensible access
summary is not utilization alone. It should report served rate, offered wait,
no-offer share, and the arrival outcome decomposition together. The current
simulation results describe the mechanics of this model under severe capacity
pressure; they should not be interpreted as calibrated clinic estimates.

\newpage

\section{Takeaway 1: Utilization Is Not Access}

The baseline demand-to-capacity ratio is the key fact: the model sends \(100\)
arrivals per day into \(32\) available daily slots. This is a deliberate stress
test. It is useful for understanding metric behavior, but it is not a calibrated
clinic setting.

\[
  \text{average utilization}
  = \frac{\text{completed visits}}{\text{available slots}},
  \qquad
  \text{served rate}
  = \frac{\text{completed visits}}{\text{arrivals}}.
\]

Because the denominators differ, utilization can remain high while access is
poor. The service system can fill most slots and still serve only a small share
of patients when demand is several times larger than supply. At baseline, the
maximum possible served rate is near \(0.32\); the observed value is \(0.269\)
after no-shows and cancellations. This is why utilization alone would be a
misleading summary of patient access in this scenario.

\begin{table}[h]
\centering
\small
\renewcommand{\arraystretch}{1.12}
\begin{tabularx}{0.92\textwidth}{Yrr}
\toprule
Metric & Baseline value & Interpretation \\
\midrule
Average utilization & 0.839 & Completed visits per available slot \\
Overall served rate & 0.269 & Completed visits per arrival \\
Mean offered wait & 9.30 days & Delay among patients receiving an offer \\
Mean accepted wait & 8.35 days & Delay among patients accepting the offer \\
Overall balking rate & 0.356 & Offered patients who reject long delays \\
Class served-rate gap & 0.001 & Near zero under symmetric class parameters \\
\bottomrule
\end{tabularx}
\caption{Baseline FCFS stress-test metrics from the main note.}
\end{table}

\section{Takeaway 2: The Main Losses Are Balking and Cancellation}

Every arrival is assigned to one outcome category. The baseline decomposition
shows that most losses occur after a slot is offered or booked. No-offer is
zero, so the baseline failure mode is not horizon saturation.

\begin{table}[h]
\centering
\small
\renewcommand{\arraystretch}{1.12}
\begin{tabularx}{0.92\textwidth}{Yrl}
\toprule
Outcome & Share of arrivals & Main interpretation \\
\midrule
Served & 0.269 & Completed visit \\
Balked & 0.356 & Received an offer, rejected the delay \\
Canceled & 0.323 & Booked, then canceled before service \\
No-show & 0.052 & Booked, did not attend \\
No-offer & 0.000 & Horizon full before any offer \\
Unresolved & \(<0.001\) & Not resolved before simulation end \\
\bottomrule
\end{tabularx}
\caption{Baseline arrival outcome decomposition.}
\end{table}

The cancellation share is high, but it is mechanically consistent with the
model's daily cancellation rule and long booking delays. With a mean accepted
wait of about \(8.35\) days and a \(10\%\) daily cancellation probability, a
rough survival calculation gives \(0.9^{8.35}\approx0.416\), or an implied
cumulative cancellation risk near \(58\%\) among long-wait bookings. The
observed cancellation share among booked patients is therefore directionally
consistent with the implementation.

The important interpretation is not that cancellation is empirically realistic.
It is that long booking delays compound daily cancellation exposure, and the
simulation is behaving consistently with that assumption.

\section{Takeaway 3: Removing Balking Does Not Recover Capacity}

The cleanest diagnostic in the note is the 30-seed comparison between the
baseline and a no-balking counterfactual. If balking were the main capacity
bottleneck, suppressing it should increase completed utilization. It does not.

\begin{table}[h]
\centering
\small
\renewcommand{\arraystretch}{1.10}
\begin{tabularx}{\textwidth}{Yrrr}
\toprule
Metric & Baseline & No-balking & Change \\
\midrule
Average utilization & 0.8418 & 0.8404 & \(-0.0014\) \\
Overall served rate & 0.2692 & 0.2682 & \(-0.0011\) \\
Booked rate & 0.6450 & 0.7034 & \(+0.0585\) \\
Balked share & 0.3550 & 0.0000 & \(-0.3550\) \\
Canceled share & 0.3248 & 0.3836 & \(+0.0588\) \\
No-offer share & 0.0000 & 0.2966 & \(+0.2966\) \\
Mean accepted wait & 8.33 d & 9.45 d & \(+1.12\) d \\
\bottomrule
\end{tabularx}
\caption{Controlled no-balking diagnostic over 30 seeds.}
\end{table}

The mechanism is a shift in accounting. Without balking, more patients accept
appointments, but those bookings extend farther into the horizon, increase
cancellation exposure, and create no-offer arrivals once the horizon fills.
Completed visits remain essentially unchanged. This is why balking should be
treated as an access-allocation and loss-channel mechanism, not as a direct
aggregate utilization lever.

\section{Takeaway 4: Driver Hierarchy and Reporting Implications}

The one-at-a-time sensitivity slices and randomized regression screen are
consistent. They identify different mechanisms rather than one universal
``best'' metric.

\begin{table}[h]
\centering
\small
\renewcommand{\arraystretch}{1.12}
\begin{tabularx}{\textwidth}{L{0.23\textwidth}L{0.31\textwidth}Y}
\toprule
Result target & Strongest evidence & Interpretation \\
\midrule
Overall served rate & Arrival rate coefficient \(\beta=-0.774\) & Demand pressure dominates access when capacity is fixed. \\
Mean offered wait & Arrival rate coefficient \(\beta=+0.576\) & More demand pushes offers farther into the booking horizon. \\
Utilization & No-show threshold coefficients \(+0.342\), \(+0.366\); no-show step coefficients \(-0.289\), \(-0.294\) & No-show behavior directly changes completed visits because missed slots are not rebooked. \\
Balking effect & Balking-step utilization effect near zero in OAT; separate regression \(\beta=+0.016\), \(p=0.72\) & Balking reallocates access and loss channels more than it changes aggregate slot completion. \\
Class access gap & Behavioral parameters have opposite-sign own-class and cross-class effects & FCFS is neutral in rule form, but shared capacity transmits behavioral differences into class-level access gaps. \\
\bottomrule
\end{tabularx}
\caption{Highest-signal sensitivity findings from the main note.}
\end{table}

Two reporting choices follow directly from these results. First, offered wait
should be preferred over accepted wait when diagnosing congestion because
offered delay is recorded before the balking decision. Accepted wait excludes
patients who rejected high-delay offers, so it can appear to improve exactly
when access is worsening. Second, offered wait should always be paired with
no-offer share. Once the horizon saturates, no-offer patients receive no delay
information and are excluded from wait metrics entirely.

\section{Caveats for Interpretation}

\begin{itemize}
  \item The baseline is a stress test, not a calibrated clinical scenario. The
  \(3.1{:}1\) demand-to-capacity ratio is intentionally severe.
  \item The findings describe the mechanics of the current FCFS formulation.
  They should not be interpreted as causal estimates for real appointment
  systems.
  \item The symmetric baseline has essentially no class gap, as expected.
  Class-gap estimates near this baseline are noisy with few seeds; stronger
  equity conclusions require more replicated simulations.
\end{itemize}

\section*{Where to Dig in the Full Note}

\begin{tabularx}{\textwidth}{L{0.33\textwidth}Y}
\toprule
Brief topic & Full-note reference \\
\midrule
Baseline setup and metric definitions & Sections 2.1--2.3; Appendix A \\
Utilization vs.\ served-rate gap & Section 2.3; Section 5 \\
Baseline outcome decomposition & Section 3.1; Table 2 \\
Balking suppression diagnostic & Section 3.2; Table 3; Section 7 for the balking deep dive \\
Offered vs.\ accepted wait & Sections 4 and 5.1--5.3 \\
Sensitivity driver hierarchy & Sections 6.1--6.5; Appendix D for the regression table \\
Class-gap interpretation & Sections 6.5 and 8 \\
Robust findings and caveats & Sections 9.1--9.2; Section 10 for discussion questions \\
\bottomrule
\end{tabularx}

\end{document}
